\section{Opérateur à deux couches et hiérarchie infinie}
\label{sec:jerarquia-espectral-infinita-rev7}
\index{tours!hiérarchie infinie}\index{semi-groupe spectral}

Soient
\[
 x=A_{\rm rad},\qquad y=C^*_{\rm rad},\qquad
 R_{\rm ch}^2=I,\qquad
 P_\pm^{\rm ch}=\frac{I\pm R_{\rm ch}}2.
\]
Le lecteur des deux feuillets est réuni dans l'opérateur
\[
 T_{\rm ang}=\exp(-xI-yR_{\rm ch})
 =q_+P_+^{\rm ch}+q_-P_-^{\rm ch},
\]
avec
\[
 q_-=e^{-x+y},\qquad q_+=e^{-x-y}.
\]
Ses moments spectraux pair et impair sont
\[
 c_n=\operatorname{Tr}(T_{\rm ang}^n)=q_-^n+q_+^n,\qquad
 s_n=-\operatorname{Tr}(R_{\rm ch}T_{\rm ang}^n)=q_-^n-q_+^n.
\]
L'involution \(C^*\mapsto-C^*\) conserve \(c_n\) et inverse \(s_n\). De manière équivalente,
\[
 c_n=2e^{-nx}\cosh(ny),\qquad
 s_n=2e^{-nx}\sinh(ny).
\]
Il s'agit de la convention fixée dans le chapitre consacré aux canaux. Pour déclarer sans homonymie la présentation alternative des sources relatives aux tours, posons
\[
 R_{\rm tow}:=-R_{\rm ch},\qquad
 P_\pm^{\rm tow}=P_\mp^{\rm ch},
\]
\[
 T_{\rm tow}:=e^{-xI+yR_{\rm tow}}=T_{\rm ang},\qquad
 \operatorname{Tr}(R_{\rm tow}T_{\rm tow}^n)
 =-\operatorname{Tr}(R_{\rm ch}T_{\rm ang}^n)=s_n.
\]
Les deux écritures décrivent le même lecteur scalaire et ne constituent pas une seconde orientation des tours.

\begin{theorem}[Addition, norme et récurrence]
\label{thm:recurrencia-bicapa-rev7}
Pour \(m,n\ge0\),
\[
 c_{m+n}=\frac{c_mc_n+s_ms_n}{2},\qquad
 s_{m+n}=\frac{s_mc_n+c_ms_n}{2},
\]
\[
 c_n^2-s_n^2=4e^{-2nx}.
\]
Si \(p=q_-q_+=e^{-2x}\), les deux suites satisfont
\[
 X_{n+2}=c_1X_{n+1}-pX_n.
\]
De plus,
\[
 \partial_Ac_n=-\frac{n\pi}{180}c_n,
 \qquad
 \partial_As_n=-\frac{n\pi}{180}s_n,
\]
\[
 \partial_{C^*}c_n=\frac{n\pi}{180}s_n,
 \qquad
 \partial_{C^*}s_n=\frac{n\pi}{180}c_n.
\]
\end{theorem}

\begin{proof}
Les identités d'addition résultent du développement de \((q_-^m\pm q_+^m)(q_-^n\pm q_+^n)\). La norme est
\[
 (q_-^n+q_+^n)^2-(q_-^n-q_+^n)^2
 =4(q_-q_+)^n.
\]
La dernière égalité est la récurrence déterminée par les racines caractéristiques \(q_-\) et \(q_+\). Les équations différentielles s'obtiennent en dérivant les expressions hyperboliques ; le facteur \(\pi/180\) appartient à la carte angulaire, et non à une sélection d'espèce.
\end{proof}

\begin{theorem}[Semi-groupe complet des hauteurs]
\label{thm:semigrupo-alturas-rev7} Les hauteurs de la hiérarchie transversale forment
\[
 \mathfrak S
 =\langle90,120\rangle
 =\{90a+120b:a,b\in\mathbb Z_{\ge0},\ a+b>0\}.
\]
De manière équivalente,
\[
 \boxed{\mathfrak S=\{90,120\}\cup\{30r:r\ge6\}.}
\]
\end{theorem}

\begin{proof}
En divisant par \(30\), on obtient le semi-groupe \(\langle3,4\rangle\). Ses éléments positifs sont \(\{3,4\}\cup\{r:r\ge6\}\) : \(5\) est l'unique lacune positive postérieure aux générateurs, et tout \(r\ge6\) s'obtient par récurrence en ajoutant \(3\) ou \(4\). En multipliant de nouveau par \(30\), on obtient la formule.
\end{proof}

La présentation conserve également la généalogie d'une hauteur :
\[
 \mathfrak S\cong
 \bigl(\mathbb Z_{\ge0}^2\setminus\{(0,0)\}\bigr)
 \big/\bigl((a+4,b)\sim(a,b+3)\bigr),
\]
puisque \(4\cdot90=3\cdot120=360\). La hauteur ne détermine pas à elle seule la composition qui l'a produite ; les généalogies \(4\cdot90\) et \(3\cdot120\) restent distinctes dans le registre chronologique enrichi, bien qu'elles donnent le même indice.

Pour expliciter la récurrence échantillonnée au pas commun de trente, définissons
\[
 a_{30}:=q_-^{30},\qquad b_{30}:=q_+^{30},\qquad
 u_r:=s_{30r}=a_{30}^r-b_{30}^r.
\]
Alors
\[
 \boxed{u_{r+2}=(a_{30}+b_{30})u_{r+1}-a_{30}b_{30}u_r.}
\]
Le terme auxiliaire \(u_5=s_{150}\) peut intervenir dans la propagation de la récurrence, mais \(150\notin\mathfrak S\) ; une identité de calcul ne crée pas de nouvelle hauteur admissible.

En particulier,
\[
90,120,180,210,240,270,300,330,360,390,420,450,480,\ldots
\]
sont les premières hauteurs. Les huit hauteurs mises en évidence dans les exposés antérieurs sont des premières lectures de \(\mathfrak S\), et non sa définition ou sa limite.

À titre de guide de lecture, sa généalogie initiale peut se lire sans être assimilée au domaine complet :
\[
\begin{array}{c@{\quad}c@{\quad}l}
90&3\cdot30&\text{générateur minimal}\\
120&4\cdot30&\text{générateur minimal}\\
180&2\cdot90&\text{premier mode pair}\\
210&90+120&\text{mode composé}\\
240&2\cdot120&\text{second mode pair}\\
270&3\cdot90&\text{répétition de la branche 90}\\
360&3\cdot120&\text{répétition de la branche 120}\\
420&2(90+120)&\text{retour composé}
\end{array}
\]
La dernière colonne consigne la provenance structurelle ; elle n'introduit pas de paramètres et n'affirme pas que ces huit lignes épuisent la hiérarchie. Les mêmes entiers \(90,120\) sont des cardinalités d'incidence dans \(\mathcal F_6\) et \(\mathcal F_8\), mais cette coïncidence n'identifie pas l'opérateur à deux couches à l'inclusion \(\iota_{6,8}\).

L'interface des deux générateurs se réduit au couple
\[
 \mathcal K_{\rm ang}^{90,120}(q_+,q_-):=(s_{90},s_{120}).
\]
Elle ne doit être confondue ni avec les résolvantes de deux sous-queues ni avec la fonctionnelle ultérieure \(K_{6\to8}=12\Delta S_{90}-\Delta S_{120}\), dont la définition et la démonstration se trouvent dans le \cref{mdg:chap:barbero}. Les sous-queues y appartiennent à la complétion analytique de la hiérarchie ; elles ne sont ni de nouveaux générateurs ni des coefficients d'espèce.

\begin{definition}[Module convergent de raffinements spectraux]
\label{def:modulo-refinamientos-torre-rev7} Le domaine intégral du relèvement global est
\[
 \mathcal E_{\mathfrak S}
 :=\left\{\eta\in\mathbb Z^{\mathfrak S}:
 \sum_{h\in\mathfrak S}|\eta_hs_h|<\infty\right\},
 \qquad s_h=q_-^h-q_+^h.
\]
Chaque \(\eta\in\mathcal E_{\mathfrak S}\) détermine le caractère spectral
\[
 \mathcal R_\eta
 :=\exp\!\left(\sum_{h\in\mathfrak S}\eta_hs_h\right).
\]
\end{definition}

La loi des masses ne réside pas dans cette définition : sa signature centrale \((n_A,s_C,k_{\Delta_4},b_{120},b_\zeta)\in\mathbb Z^5\) est construite auparavant à partir du registre. Le relèvement par \(\eta\) agit ensuite et conserve la provenance du terme de base \(b_{120}s_{120}\). En particulier, \(N_{120}=b_{120}+\eta_{120}\) est le coefficient total évalué, et non une nouvelle coordonnée de signature ou une décomposition inverse unique.

La complétude du semi-groupe ne sélectionne pas à elle seule une espèce. Lorsque \(\eta\) est obtenu à partir d'un résidu métrologique fourni, il s'agit d'une reconstruction calibrée exacte ; la portée prospective exige que la route, la retenue et l'enroulement produisent \(\eta\) avant l'examen de ce résidu. La récurrence précédente engendre tous les moments une fois le raffinement fixé, sans introduire une loi distincte à chaque hauteur.
